\documentclass[11pt]{article}

\usepackage[T1]{fontenc}
\usepackage[utf8]{inputenc}
\usepackage{lmodern}
\usepackage{microtype}
\usepackage[margin=1in]{geometry}
\usepackage{amsmath,amssymb,amsthm,mathtools}
\usepackage{booktabs}
\usepackage{longtable}
\usepackage{array}
\usepackage{enumitem}
\usepackage{xcolor}
\usepackage{listings}
\usepackage{fancyhdr}
\usepackage{lastpage}
\usepackage{hyperref}
\usepackage[nameinlink,noabbrev]{cleveref}
\usepackage{url}

\hypersetup{
  colorlinks=true,
  linkcolor=blue!55!black,
  citecolor=blue!55!black,
  urlcolor=blue!55!black,
  pdftitle={Ramified Dynamic Operator Reflection: A Proposed Contender 14 for BusyLean},
  pdfsubject={Lean 4 large-number construction and formal comparison},
  pdfkeywords={Lean 4, BusyLean, large numbers, reflection, System F omega, proof theory}
}

\pagestyle{fancy}
\fancyhf{}
\lhead{BusyLean Contender 14}
\rhead{Ramified dynamic operator reflection}
\cfoot{\thepage\ of \pageref{LastPage}}
\setlength{\headheight}{14pt}

\newtheorem{theorem}{Theorem}[section]
\newtheorem{lemma}[theorem]{Lemma}
\newtheorem{proposition}[theorem]{Proposition}
\newtheorem{corollary}[theorem]{Corollary}
\theoremstyle{definition}
\newtheorem{definition}[theorem]{Definition}
\newtheorem{remark}[theorem]{Remark}

\newcommand{\Nat}{\ensuremath{\mathbb{N}}}
\newcommand{\Family}{\ensuremath{\mathsf{Family}}}
\newcommand{\Transformer}{\ensuremath{\mathsf{Transformer}}}
\newcommand{\Prog}{\ensuremath{\mathsf{Progressive}}}
\newcommand{\Fd}{\ensuremath{\mathsf{F}_{\mathrm{dyn}}}}
\newcommand{\lean}[1]{\texttt{#1}}
\newcommand{\Cthirteen}{\ensuremath{C_{13}}}
\newcommand{\Cfourteen}{\ensuremath{C_{14}}}
\newcommand{\Fthirteen}{\ensuremath{F_{13}}}
\newcommand{\Ffourteen}{\ensuremath{F_{14}}}
\newcommand{\opstage}{\ensuremath{\mathsf{operatorStage}}}
\newcommand{\pack}{\ensuremath{\mathsf{packFamily}}}
\newcommand{\prefixmax}{\ensuremath{\mathsf{prefixMax}}}
\newcommand{\tower}{\ensuremath{\mathsf{operatorTower}}}
\newcommand{\step}{\ensuremath{\mathsf{towerStep}}}

\lstdefinelanguage{Lean}{
  morekeywords={def,theorem,lemma,namespace,end,import,open,abbrev,inductive,structure,
    where,match,with,if,then,else,fun,by,exact,have,show,change,rw,calc,forall},
  sensitive=true,
  morecomment=[l]{--},
  morecomment=[s]{/-}{-/},
  morestring=[b]"
}
\lstset{
  language=Lean,
  basicstyle=\ttfamily\small,
  keywordstyle=\bfseries\color{blue!50!black},
  commentstyle=\itshape\color{green!35!black},
  stringstyle=\color{red!45!black},
  columns=fullflexible,
  keepspaces=true,
  breaklines=true,
  frame=single,
  framerule=0.3pt,
  xleftmargin=0.4em,
  xrightmargin=0.4em,
  aboveskip=0.8em,
  belowskip=0.8em
}

\title{\textbf{Ramified Dynamic Operator Reflection}\\
A Proposed Contender 14 for BusyLean}
\author{Research note prepared for the BusyLean project}
\date{12 October 2026\\\small Base repository commit \texttt{1ee50c7d97bbed57a60709f661453c3e213f4db3}}

\begin{document}
\maketitle

\begin{abstract}
We propose a closed computable Lean natural number \lean{contender\_14} that
strictly exceeds the current BusyLean champion, \lean{contender\_13}, without
using the predecessor constant in the new definition.  The construction is a
ramified tower of operator-level reflection over the current dynamic-kind
System~$F_\omega$ normalizer.  Its base step upgrades the unary oracle of
Contender~13 to a packed countable family and explicitly dominates every
visible family coordinate.  A strict meta-step then feeds the transformer's
whole output family into a full operator-stage closure along a schedule stacked
above Contender~12's autonomous schedule.  Repeating this meta-step $n+1$ times
at input $n$ yields a strict tower of progressively stronger transformers.
The seed family exposes all finite additional iterations of the Contender~13
construction.

The proposed Lean development proves, without evaluating any contender,
\[
  \lean{contender\_13} < \lean{contender\_14}.
\]
It also proves the much sharper certificate
\[
  \Fthirteen^{\,42}(\Cthirteen) < \Ffourteen(42) < \Cfourteen,
\]
where the first exponent denotes forty-two right-associated functional
iterations.  Moreover, every accepted dynamic-kind program of depth at most
$n$, run with the packed iterated seed family as oracle, has value below
$\Ffourteen(n)$.  All numerical definitions are structural and executable in
principle.  The archive contains Lean source, integration scripts, an axiom and
dependency audit extension, and this article.  The PDF was compiled and
inspected.  The preparation environment did not contain Lean~4.34.0, so no
kernel-build receipt is fabricated; final acceptance requires running the
included commands in the pinned BusyLean checkout.
\end{abstract}

\tableofcontents
\newpage

%======================================================================
\section{Executive statement}

The proposed number is
\begin{lstlisting}
def contender_14 : Nat := F14 (F14 42)
\end{lstlisting}
with the following constructional ingredients:
\begin{lstlisting}
def family13 (n : Nat) : Nat := F13 (F13 n)

def iterateF13 : Nat -> Nat -> Nat
  | 0, n     => n
  | k + 1, n => F13 (iterateF13 k n)

def seed13 : W.Family :=
  fun i n => iterateF13 i (family13 n)

def operatorTower : Nat -> Transformer
  | 0     => fdFamilyTransformer
  | k + 1 => towerStep (operatorTower k)

def F14 (n : Nat) : Nat :=
  operatorTower (n + 1) seed13 n n
\end{lstlisting}
The base transformer is a family-aware version of the repository's
\lean{Fd.reflectFd}.  The schedule \lean{tree14} is obtained by placing another
42-level autonomous hyper-closure above \lean{tree12}.  The successor operation
\lean{towerStep} does not merely compose the previous transformer with a fast
arithmetic function; it reruns the repository's full operator-level reflection
machinery with the previous transformer's output as the new oracle family.

The main theorem is represented by the Lean declaration
\begin{lstlisting}
theorem contender_13_lt_contender_14 :
    contender_13 < contender_14
\end{lstlisting}
and follows from three uniform facts:
\begin{enumerate}[label=(\roman*)]
  \item every level of \lean{operatorTower} is progressive;
  \item the seed's coordinate zero is \lean{family13}; and
  \item \lean{F14} is strictly input-increasing.
\end{enumerate}
No numeral is evaluated in this proof.

%======================================================================
\section{Repository context and admissibility constraints}

BusyLean is a Lean~4 continuation of the ``name the biggest number'' contest.
Its operative rules require a closed computable \lean{Nat}, a formal strict
comparison with the previous contender, a small permitted axiom set, and a
source-level prohibition on defining a new entry from the predecessor constant
itself.  The repository currently identifies Contender~13 as the champion.
That contender reflects all bounded programs of dynamic-kind System~$F_\omega$
over the function \lean{family12}; it is therefore already far beyond an
ordinary tower, Ackermann function, or fixed fast-growing hierarchy at a
hand-written ordinal notation.

The relevant predecessor definitions are schematically
\begin{align*}
  \mathsf{family}_{12}(n) &= F_{12}(F_{12}(n)),\\
  F_{13}(n) &= \mathsf{reflectFd}(\mathsf{family}_{12},n),\\
  C_{13} &= F_{13}(F_{13}(42)).
\end{align*}
The reflector contains a finite enumeration of terms, a sound bounded type
checker, a total normalizer justified by strong normalization, and a maximum of
the decoded numeral values.  Thus any serious successor should operate on the
\emph{construction} of $F_{13}$, not append an elementary wrapper to $C_{13}$.

\subsection{Design requirements used here}

The proposal was designed against the following stronger checklist.

\begin{longtable}{>{\raggedright\arraybackslash}p{0.24\textwidth}
                  >{\raggedright\arraybackslash}p{0.68\textwidth}}
\toprule
Requirement & Realization \\
\midrule
\endhead
Closed and constructive & Every list, maximum, normalization, tree recursion,
and numerical iteration is finite or structurally recursive. \\
No predecessor constant & \lean{Contender14.lean} uses \lean{F13} and
\lean{family12}'s construction, but not \lean{contender\_13}. \\
Formal margin & The proof gives both $C_{13}<C_{14}$ and a 42-extra-iteration
lower bound. \\
Qualitative amplification & Reflection is upgraded from a unary oracle to a
family oracle, lifted through operator reflection, and nested by an
input-indexed tower of strict meta-steps. \\
Auditability & Definitions and comparison proofs are separated; the archive
adds axiom checks and a transitive dependency-closure check. \\
Budget plausibility & Earlier large modules are imported precompiled.  New
definitions are short transparent recursors; proofs reuse generic progressivity
theorems. \\
No informal oracle & The packed family is an ordinary computable function, and
program evaluation is the existing total \Fd\ evaluator. \\
\bottomrule
\end{longtable}

%======================================================================
\section{Why obvious alternatives are inadequate}

\subsection{Arithmetic post-composition}
A definition such as $2^{C_{13}}$, a tetration of height $C_{13}$, or even an
Ackermann diagonal applied to $C_{13}$ would certainly be larger if it were
allowed.  It would nevertheless violate the contest's no-laziness principle:
the new entry would obtain its strength from the already named constant rather
than from a stronger constructive mechanism.  Replacing the constant by the
definitional equal term $F_{13}(F_{13}(42))$ would evade only a textual check,
not the mathematical objection.

\subsection{Unbounded Busy Beaver maxima}
A true maximum over all terminating programs below a syntax bound is not
computable because termination is undecidable.  The repository avoids this by
using total typed languages and formal normalization theorems.  The present
proposal keeps that discipline: it reuses the total \Fd\ evaluator and only
maximizes over a finite syntax list.

\subsection{A new calculus without a normalization proof}
A hierarchy of dynamic universes $\mathtt{dyn}_0,\mathtt{dyn}_1,\ldots$ is a
promising qualitative successor and is discussed in the repository's roadmap.
However, introducing it responsibly requires extending syntax, kinding,
conversion, type inference, semantic interpretation, and the reducibility
argument.  A large numeral whose central language theorem is merely asserted
would not be a valid contest entry.  The current proposal instead extracts a
substantial new reflection level from machinery that is already normalized and
audited, while leaving the dynamic-universe ladder as the leading next base-
language project.

%======================================================================
\section{Preliminaries: families, transformers, and progressivity}

A family is a two-argument natural-number function
\[
  \Family := \Nat \to \Nat \to \Nat.
\]
The first argument is an oracle-family coordinate and the second is the
numerical input.  A transformer maps families to families:
\[
  \Transformer := \Family \to \Family.
\]
The repository's key abstract invariant is progressivity.

\begin{definition}[Progressive transformer]
A transformer $T$ is progressive when
\[
  \Prog(T) :\Longleftrightarrow
  \forall S,i,o,n\,\bigl(i\le n \Rightarrow S(i,n)<T(S)(o,n)\bigr).
\]
The output coordinate $o$ is unrestricted.  Only the seed coordinate $i$ must
be visible at the current budget.
\end{definition}

This deliberately avoids monotonicity in the oracle.  Oracle programs may
inspect values and branch, so replacing an oracle pointwise by a larger one
need not increase a program's output.  Progressivity is instead proved from an
\emph{explicit occurrence}: a visible seed value is inserted into a finite
maximum, or a child value is directly exposed to the parent transformer.

The existing \lean{operatorStage} construction is generic in a progressive
base transformer.  It first iterates a transformer along the universal
well-founded tree \lean{index}; it then recursively splices child transformers
and lifts the splice again.  The repository proves
\begin{equation}
  \Prog(T) \quad\Longrightarrow\quad
  \Prog(\opstage(T,a))
  \label{eq:operator-progressive}
\end{equation}
for every schedule tree $a$.

%======================================================================
\section{Family-aware dynamic-kind reflection}

Contender~13's \lean{Fd.reflectFd} accepts one unary oracle $f:\Nat\to\Nat$.
For operator reflection, a countable family is more natural: child stages must
remain separately addressable.  The repository already provides a computable
pairing code \lean{packFamily} satisfying the intended decode equation for
family coordinates and arguments.  We therefore define
\begin{align}
R_{\mathrm{fam}}(S,n)
  := {}& n+1+\max\Bigl(
    \max_{i\le n} S(i,n),\notag\\
   &\hspace{33mm}
    \max_{t\in\mathsf{ctmsUpTo}(n)}
      \mathsf{fvalue}(\pack(S),n,t)
  \Bigr).
\label{eq:family-reflector}
\end{align}
In Lean this is \lean{reflectFdFamily}:
\begin{lstlisting}
def reflectFdFamily (fs : W.Family) (n : Nat) : Nat :=
  n + 1 + BigNumber.maxN
    (W.prefixMax (fun i => fs i n) n)
    (BigNumber.maximum
      (Fd.fvalue (packFamily fs) n) (Fd.ctmsUpTo n))
\end{lstlisting}

\begin{lemma}[Visible-coordinate domination]
For every family $S$ and $i\le n$,
\[
  S(i,n)<R_{\mathrm{fam}}(S,n).
\]
\end{lemma}
\begin{proof}
The prefix maximum contains $S(i,n)$.  Hence $S(i,n)$ is at most the left
component of the outer maximum.  The summand $n+1$ makes the final inequality
strict.
\end{proof}

\begin{lemma}[Program domination]
If $t\in\mathsf{ctmsUpTo}(n)$, then
\[
  \mathsf{fvalue}(\pack(S),n,t)<R_{\mathrm{fam}}(S,n).
\]
\end{lemma}
\begin{proof}
The value is at most the finite maximum over the enumerated terms, which is at
most the outer maximum.  Again $n+1$ supplies strictness.
\end{proof}

The associated transformer ignores its output coordinate:
\[
  T_0(S)(o,n):=R_{\mathrm{fam}}(S,n).
\]
Visible-coordinate domination is exactly \(\Prog(T_0)\).  Thus
\lean{fdFamilyTransformer} can be fed to all generic operator-reflection
lemmas without changing the existing \Fd\ normalizer.

%======================================================================
\section{A schedule above Contender 12's schedule}

Contender~12 uses
\[
  \mathsf{tree}_{12}
  =\mathsf{realize}\bigl(\mathsf{hyper}(42,
       \mathsf{encode}(\mathsf{index}))\bigr).
\]
The autonomous notation language contains countably branching trees,
syntactic addition and multiplication, $\omega$-powers, derivatives, a
finite critical hierarchy, a common-critical-point closure, and higher
derivatives.  It is intentionally used as tree syntax; no quotient-level
ordinal identity is needed for correctness.

The new schedule stacks the same closure once more:
\[
  \mathsf{tree}_{14}
  =\mathsf{realize}\bigl(\mathsf{hyper}(42,
       \mathsf{encode}(\mathsf{tree}_{12}))\bigr).
\]
The definition is transparent and structurally recursive.  Its role is not in
the elementary strict comparison---progressivity holds for every schedule---but
in making each meta-step a much richer closure than one fixed lift.

\begin{remark}
The exact proof-theoretic ordinal denoted by these raw trees is deliberately not
asserted.  The construction is valid because \lean{operatorStage} recurses on
the actual well-founded Lean tree, not because an informal ordinal label has
been proved correct.
\end{remark}

%======================================================================
\section{The strict meta-step}

For a progressive transformer $T$, define
\begin{equation}
  \step(T)(S)
    := \opstage(T,\mathsf{tree}_{14})\bigl(T(S)\bigr).
\label{eq:tower-step}
\end{equation}
This is \lean{towerStep}.  It has two conceptually distinct parts:
\begin{enumerate}[label=(\alph*)]
  \item first expose the entire output family $T(S)$;
  \item then use that family as the seed of the operator-stage closure of the
        same transformer $T$.
\end{enumerate}
The feedback is what gives a clean strict chain.

\begin{theorem}[Strict domination by one meta-step]
If \(\Prog(T)\), $i\le n$, and $S$ is any family, then
\[
  T(S)(i,n)<\step(T)(S)(o,n)
\]
for every output coordinate $o$.
\end{theorem}
\begin{proof}
By \cref{eq:operator-progressive}, the transformer
$\opstage(T,\mathsf{tree}_{14})$ is progressive.  Apply its progressivity to
the seed family $T(S)$ and visible coordinate $i$.
\end{proof}

\begin{theorem}[Preservation of progressivity]
If \(\Prog(T)\), then \(\Prog(\step(T))\).
\end{theorem}
\begin{proof}
For $i\le n$,
\[
  S(i,n)<T(S)(i,n)<\step(T)(S)(o,n).
\]
The first inequality is progressivity of $T$; the second is the preceding
theorem.
\end{proof}

This proof is important methodologically.  It never assumes that reflection is
monotone in its oracle.  Both strict inequalities are witnessed by the generic
progressivity interface.

%======================================================================
\section{The ramified operator tower}

Define a sequence of transformers by
\begin{align*}
  T_0 &:= \mathsf{fdFamilyTransformer},\\
  T_{k+1} &:= \step(T_k).
\end{align*}
This is \lean{operatorTower}.  Structural recursion is on $k$.

\begin{theorem}[Tower invariant]
For every $k\in\Nat$, $T_k$ is progressive.
\end{theorem}
\begin{proof}
Induct on $k$.  The base case is visible-coordinate domination for
\lean{reflectFdFamily}.  The successor case is preservation of progressivity
by \lean{towerStep}.
\end{proof}

\begin{theorem}[Strict level chain]
For every $k$, family $S$, $i\le n$, and output coordinate $o$,
\[
  T_k(S)(i,n)<T_{k+1}(S)(o,n).
\]
\end{theorem}
\begin{proof}
This is strict domination by one meta-step, using the tower invariant at level
$k$.
\end{proof}

The word ``tower'' here is therefore not rhetorical.  The Lean development
contains a pointwise strict theorem between every two successive levels.  Each
successor is also substantially more structured than a simple function
composition: it invokes a complete operator-stage closure whose base is the
entire preceding level.

%======================================================================
\section{The iterated Contender 13 seed}

The predecessor construction is exposed as
\[
  \mathsf{family}_{13}(n):=F_{13}(F_{13}(n)).
\]
Thus $\mathsf{family}_{13}(42)=C_{13}$ definitionally, but the constant
\lean{contender\_13} is absent from the new executable definition.

Define right-associated iteration by
\begin{align*}
  \mathsf{iterateF13}(0,x)&:=x,\\
  \mathsf{iterateF13}(k+1,x)&:=F_{13}(\mathsf{iterateF13}(k,x)).
\end{align*}
The seed family is
\begin{equation}
  S_{13}(i,n):=
  \mathsf{iterateF13}\bigl(i,\mathsf{family}_{13}(n)\bigr).
\label{eq:seed13}
\end{equation}
Coordinate zero is exactly the predecessor family.  Coordinate $i$ supplies
$i$ further applications of the dynamic-kind reflector.  Because visible
coordinates satisfy $i\le n$, the input $n$ makes the first $n+1$ such
iterations simultaneously available to the packed oracle and to the
progressivity maximum.

\begin{lemma}[Strictness of nonempty \(F_{13}\)-iteration]
For all $k,x\in\Nat$,
\[
  x<\mathsf{iterateF13}(k+1,x).
\]
\end{lemma}
\begin{proof}
The repository proves $y<F_{13}(y)$ for every $y$.  Induct on $k$ and append
this input-growth inequality to the preceding chain.  No monotonicity of
$F_{13}$ is required.
\end{proof}

%======================================================================
\section{\texorpdfstring{Definition of \(F_{14}\) and the final number}{Definition of F14 and the final number}}

At input $n$, choose tower level $n+1$, seed it with $S_{13}$, select output
coordinate $n$, and evaluate at numerical input $n$:
\begin{equation}
  F_{14}(n):=T_{n+1}(S_{13})(n,n).
\label{eq:F14}
\end{equation}
Finally set
\begin{equation}
  C_{14}:=F_{14}(F_{14}(42)).
\label{eq:C14}
\end{equation}
The outer self-application mirrors the recent BusyLean contenders and is useful
because $F_{14}$ is formally input-increasing.

The complete executable spine is only a few transparent definitions.  The
size comes from the already verified total interpreters and from structural
reflection, not from a giant literal term.

%======================================================================
\section{Formal comparison theorem}

\begin{theorem}[Visible seed bound]
If $i\le n$, then
\[
  S_{13}(i,n)<F_{14}(n).
\]
\end{theorem}
\begin{proof}
Apply progressivity of $T_{n+1}$ to seed $S_{13}$, input coordinate $i$, output
coordinate $n$, and numerical input $n$.
\end{proof}

At coordinate zero this gives
\begin{equation}
  \mathsf{family}_{13}(n)<F_{14}(n).
\label{eq:family13-F14}
\end{equation}
The existing input-growth theorem for $F_{13}$ yields
\[
  n<F_{13}(n)<F_{13}(F_{13}(n))
    =\mathsf{family}_{13}(n).
\]
Combining with \cref{eq:family13-F14} proves
\begin{equation}
  n<F_{14}(n).
\label{eq:F14-input}
\end{equation}

\begin{theorem}[Main comparison]
\[
  C_{13}<C_{14}.
\]
\end{theorem}
\begin{proof}
At $n=42$, \cref{eq:family13-F14} gives
\[
  C_{13}=\mathsf{family}_{13}(42)<F_{14}(42).
\]
Apply \cref{eq:F14-input} at input $F_{14}(42)$:
\[
  F_{14}(42)<F_{14}(F_{14}(42))=C_{14}.
\]
Transitivity completes the proof.
\end{proof}

The corresponding Lean term is intentionally short:
\begin{lstlisting}
theorem contender_13_lt_contender_14 :
    contender_13 < contender_14 :=
  Nat.lt_trans (family13_lt_F14 42)
    (input_lt_F14 (F14 42))
\end{lstlisting}
The complexity lies in the reusable invariants, not in forcing Lean to
normalize either number.

%======================================================================
\section{A certified large margin}

The coordinate $i=42$ is visible at input $42$.  Therefore the visible-seed
bound gives the exact theorem
\begin{equation}
  \mathsf{iterateF13}(42,C_{13})<F_{14}(42).
\label{eq:42-margin}
\end{equation}
Since $C_{13}=F_{13}^2(42)$, the left side is $F_{13}^{44}(42)$ under the
right-associated iteration convention.  Thus the formal chain includes
\[
  C_{13}<F_{13}^{44}(42)<F_{14}(42)<C_{14}.
\]
The archive exposes \cref{eq:42-margin} as
\lean{fortyTwo\_F13\_steps\_lt\_F14}.

This lower bound is useful for two reasons.  First, it demonstrates a margin
that cannot be dismissed as ``plus one.''  Second, it is independent of any
informal ordinal assignment to dynamic-kind System~$F_\omega$ or to the
schedule trees.  It follows only from explicit occurrences in finite maxima
and from structural progressivity.

%======================================================================
\section{Program domination}

The base family reflector dominates every accepted \Fd\ term in the finite
enumeration.  The strict tower-level chain places the base below every positive
operator-tower level.  Consequently:

\begin{theorem}[Bounded program domination]
If $t\in\mathsf{ctmsUpTo}(n)$, then
\[
  \mathsf{fvalue}(\pack(S_{13}),n,t)<F_{14}(n).
\]
\end{theorem}
\begin{proof}
The finite maximum in \cref{eq:family-reflector} first gives
\[
  \mathsf{fvalue}(\pack(S_{13}),n,t)
  <T_0(S_{13})(0,n).
\]
The strict tower chain, iterated through the positive level $n+1$, gives
\[
  T_0(S_{13})(0,n)<T_{n+1}(S_{13})(n,n)=F_{14}(n).
\]
\end{proof}

At the outer budget $F_{14}(42)$ this becomes
\begin{align*}
 &t\in\mathsf{ctmsUpTo}(F_{14}(42))\\
 &\qquad\Longrightarrow\quad
 \mathsf{fvalue}(\pack(S_{13}),F_{14}(42),t)<C_{14}.
\end{align*}
This is formalized as \lean{fvalue\_lt\_contender\_14}.

The oracle is more expressive than the unary predecessor oracle.  A program
can decode an index and query any seed coordinate, while operator-stage nodes
replace the seed by families of distinct child outputs.  The family interface
therefore preserves dynamic access rather than collapsing all child stages to
one diagonal function.

%======================================================================
\section{Constructivity and computability}

\begin{proposition}
\lean{contender\_14} denotes a total computable natural number in principle.
\end{proposition}
\begin{proof}[Computational analysis]
Each component is total:
\begin{enumerate}[label=(\arabic*)]
  \item \lean{Fd.ctmsUpTo n} is a finite list produced by recursion on $n$.
  \item The bounded checker and the \Fd\ evaluator are total; accepted terms
        normalize by the repository's strong-normalization theorem.
  \item \lean{BigNumber.maximum} and \lean{W.prefixMax} inspect finite data.
  \item \lean{runTransformer}, \lean{operatorStage}, \lean{Autonomous.encode},
        and \lean{Autonomous.realize} recurse structurally on well-founded
        inductive trees.
  \item \lean{iterateF13} and \lean{operatorTower} recurse on a natural number.
  \item Function packing is an explicit arithmetic encoding, not an oracle.
\end{enumerate}
No definition is partial, unsafe, noncomputable, opaque, externally
implemented, or based on unbounded search.  Therefore Lean's evaluator has a
finite reduction sequence for the closed term, although its resource demand is
astronomical.
\end{proof}

\begin{remark}
``Computable in principle'' does not mean feasible to print.  BusyLean's own
README warns never to evaluate the contenders.  Kernel checking verifies the
transparent definitions and proof terms without normalizing the final numeral.
\end{remark}

%======================================================================
\section{Why the growth jump is substantial}

Only the exact inequalities above are claimed as theorems.  Nevertheless, the
architecture supports a useful qualitative analysis.

\subsection{Four independent amplification axes}

\begin{enumerate}[label=\textbf{A\arabic*.},leftmargin=3.2em]
  \item \textbf{Oracle arity.}  One unary predecessor oracle is replaced by a
        packed countable family with explicit prefix domination.
  \item \textbf{Autonomous reflection.}  A family-aware \Fd\ step is iterated
        along the W-language's universal index inside each lift.
  \item \textbf{Operator reflection.}  Child transformers remain separately
        selectable and are recursively spliced along \lean{tree14}.
  \item \textbf{Ramification.}  The entire operator closure becomes the base
        of the next level, and the number of levels is $n+1$ at input $n$.
\end{enumerate}

The predecessor uses one \Fd\ reflection over \lean{family12}.  In contrast,
$F_{14}(42)$ uses level 43 of a strict operator tower, with a seed already
containing 42 extra $F_{13}$ iterations at visible coordinate 42.  The exact
lower bound captures only a small, easy-to-state fragment of the full
operator-stage computation.

\subsection{Schematic ordinal reading}
Let $\Lambda_{\Fd}$ denote, purely schematically, the closure strength of one
family-aware dynamic-kind reflection step, and let $\alpha$ and $\gamma$ denote
the heights represented by \lean{index} and \lean{tree14}.  One may view a
lift as iterating a $\Lambda_{\Fd}$-strength operation along $\alpha$, while an
operator stage recursively lifts spliced operators along $\gamma$.  Replacing
the base transformer by the resulting feedback closure and repeating produces
a hierarchy
\[
  T_0,\quad \step(T_0),\quad \step^2(T_0),\quad\ldots
\]
whose 43rd level is selected already at input 42.

This paragraph is an explanatory model, not a formal ordinal calibration.  A
sound calibration would require assigning ordinal semantics to the actual
raw-tree operations and proving preservation theorems.  The formal Lean
comparison deliberately does not depend on such an assignment.

%======================================================================
\section{Axiom and dependency discipline}

The archive extends the repository's audit in two ways.

\subsection{Axiom reports}
The proposed audit asks Lean to print the axioms of
\begin{itemize}
  \item \lean{contender\_14};
  \item \lean{contender\_13\_lt\_contender\_14};
  \item \lean{fortyTwo\_F13\_steps\_lt\_F14}; and
  \item \lean{fvalue\_lt\_contender\_14}.
\end{itemize}
No new axiom is introduced by the source.  The expected inherited set is the
repository-permitted subset of \lean{propext} and \lean{Quot.sound}.

\subsection{No-laziness dependency closure}
The existing meta-level command \lean{\#assert\_not\_in\_closure} traverses all
constants reachable through declaration types and bodies.  The integration
appendix checks that \lean{contender\_14} does not depend on any predecessor
constant \lean{contender\_13},\ldots,\lean{contender\_5}.  It is allowed to
depend on \lean{F13}, \lean{family12}, the normalizer, and the generic
reflection machinery: those are predecessor \emph{constructions}, exactly as
the contest rule permits.

\subsection{Static preflight}
The supplied Python preflight strips nested Lean comments and strings, then
checks the executable source for forbidden tokens and predecessor constants.
It also verifies the expected declaration set and records SHA-256 hashes.  This
is useful before integration but is explicitly labeled as not being a kernel
check.

%======================================================================
\section{Expected build behavior}

The repository permits earlier contender modules to be precompiled when timing
a new entry.  Under that convention the definition phase imports
\lean{Contender13}, adds one finite maximum, several small transparent
recursors, and no new normalization proof.  The comparison phase is dominated
by generic theorem instantiation and short inductions over natural tower
levels.  Thus the design is intentionally conservative with respect to the
15-second definition and 60-second proof budgets.

The archive pins the same toolchain as the base repository:
\[
  \lean{leanprover/lean4:v4.34.0}.
\]
The intended commands after applying the overlay are
\begin{lstlisting}[language={}]
python3 scripts/check.py
lake env lean test/Contender14Audit.lean
\end{lstlisting}
A successful run should be retained as a build receipt.  Because the execution
environment used to create this archive had no Lean binary and could not fetch
the 4.34.0 toolchain, this article does not claim a completed kernel build.
That limitation is explicit rather than hidden behind a fabricated log.

%======================================================================
\section{Integration plan}

The archive is an overlay, not a fork.  It contains
\begin{itemize}
  \item \lean{BusyLean/Contender14.lean};
  \item \lean{BusyLean/Comparison14.lean};
  \item \lean{test/Contender14Audit.lean};
  \item an append block for the repository-wide audit;
  \item scripts to copy the overlay and run checks;
  \item a unified integration patch and documentation;
  \item this \LaTeX{} article and compiled PDF.
\end{itemize}

The root module should import the two new modules after
\lean{BusyLean.Comparison13}.  The timing script should set
\begin{lstlisting}[language=Python]
PREVIOUS = "BusyLean.Comparison13"
CURRENT_DEF = "BusyLean.Contender14"
\end{lstlisting}
and the audit should add the declarations listed in the preceding section.

%======================================================================
\section{Risk assessment}

\begin{longtable}{>{\raggedright\arraybackslash}p{0.27\textwidth}
                  >{\raggedright\arraybackslash}p{0.28\textwidth}
                  >{\raggedright\arraybackslash}p{0.36\textwidth}}
\toprule
Risk & Consequence & Mitigation \\
\midrule
\endhead
Lean elaboration mismatch & Source needs a local syntactic repair. & The code
copies theorem shapes already used by Comparisons 11--13 and minimizes tactic
use.  Run the pinned checker before acceptance. \\
Unexpected axiom dependency & Contest audit fails. & No new semantic theorem is
introduced; all proofs use existing progressivity and elementary \lean{Nat}
lemmas.  Explicit \lean{\#print axioms} entries are supplied. \\
Definition budget regression & Entry violates timing rule. & No new object
language or normalization proof; previous modules are precompiled by the
contest's timing convention. \\
No-laziness audit failure & Entry is inadmissible. & The executable module does
not mention predecessor constants.  A transitive closure assertion is supplied. \\
Overstated proof-theoretic size & Misleading research claim. & Only strict
numeric and program bounds are theorem claims.  Ordinal discussion is labeled
schematic. \\
Oracle nonmonotonicity & Invalid comparison proof. & All inequalities use
explicit progressivity; no oracle monotonicity is assumed. \\
\bottomrule
\end{longtable}

%======================================================================
\section{Alternative constructions considered}

\subsection{\texorpdfstring{One operator stage over \Fd}{One operator stage over Fd}}
The smallest natural successor would instantiate Contender~12's generic
operator stage with family-aware \Fd\ reflection and seed it with
\lean{family13}.  That already gives a large strict improvement.  The present
proposal goes further by proving a strict feedback step and iterating that
whole step through an input-indexed tower.

\subsection{Fixed repeated reflection}
One could iterate \lean{Fd.reflectFd} 42 times.  This is simple and formally
large, but it collapses all child stages to a unary chain and fails to exploit
the repository's family and operator infrastructure.  It also gives a much
weaker structural explanation of the growth margin.

\subsection{A dynamic-kind universe ladder}
A hierarchy $\mathtt{dyn}_j$ would be a stronger base-language advance.  It is
likely the most important next long-term direction.  It is not used here
because a valid implementation requires a new normalization model and sound
checker, not just syntax.  The current operator tower can later be re-based on
such a language: replace \lean{fdFamilyTransformer} with the progressive
reflector of the verified ladder while keeping the rest of the construction.

\subsection{Large elimination and tree-returning programs}
A calculus with inductive trees and large elimination could internalize the
world/model bootstraps explored by earlier proposals.  This may eventually
outgrow any amplifier over \Fd.  It is again a base-language research program,
not a small safe overlay.

%======================================================================
\section{Further research questions}

\begin{enumerate}[leftmargin=2.1em]
  \item \textbf{Kernel-check and benchmark the overlay.}  Resolve any version-
  specific elaboration issue, record clean definition/proof times, and preserve
  complete axiom output.

  \item \textbf{Formalize a generic strict meta-closure theorem.}  Abstract
  \lean{towerStep} away from \lean{tree14} and prove reusable results for any
  schedule and progressive transformer, including closed forms for finite
  level chains.

  \item \textbf{Prove stronger program bounds at every tower level.}  Track
  not merely the base \Fd\ program maximum but programs relative to each
  derived child family, with explicit embedding lemmas through
  \lean{operatorStage}.

  \item \textbf{Develop ordinal semantics for the raw autonomous trees.}
  Define a quotient or simulation relation, prove monotonicity of \lean{add},
  \lean{mul}, \lean{pow}, \lean{derivative}, and \lean{hyper}, and derive
  certified lower bounds for \lean{tree12} and \lean{tree14}.

  \item \textbf{Implement the dynamic-universe ladder.}  Add
  $\mathtt{dyn}_0,\ldots,\mathtt{dyn}_{n-1}$ at syntax budget $n$, define the
  level-indexed semantic dependent sums, and extend strong normalization and
  inference soundness.  A diagonal reflector over this language should be the
  next qualitative base jump.

  \item \textbf{Universe-polymorphic compression.}  Investigate whether the
  finite ladder can be represented by one level-indexed kind constructor while
  preserving a decidable, predicative semantic model and fast elaboration.

  \item \textbf{Tree-capable dynamic calculus.}  Combine dynamic kinds with
  W-types or strictly positive inductives, prove normalization with large
  elimination where sound, and expose generated trees to the reflection
  schedule.

  \item \textbf{Higher meta-types.}  Formalize certified progressive
  transformers as a subtype and iterate the lift at the level
  $(\Transformer\to\Transformer)$, then at higher finite meta-types.  Compare
  this genuine type-level tower with the present same-type feedback tower.

  \item \textbf{Minimality and redundancy.}  Determine which components of
  \lean{tree14}, the seed iteration, and the feedback composition are needed
  for the strongest formal lower bounds under the timing budget.

  \item \textbf{Machine-checked no-laziness specification.}  Replace the
  current blacklist with a positive provenance policy describing which
  predecessor families and constructors may occur, and prove the policy stable
  under refactoring.

  \item \textbf{Independent normalization replay.}  Export the \Fd\ strong-
  normalization dependency graph and replay its critical semantic lemmas in a
  second small checker or a minimized Lean environment.

  \item \textbf{Compare amplifier versus language jumps.}  Build a formal
  embedding framework that can prove when an object-language extension
  dominates any bounded number of reflection amplifiers over the old language.
\end{enumerate}

%======================================================================
\section{Conclusion}

The proposed Contender~14 is a direct, auditable successor to the current
BusyLean architecture.  It avoids the prohibited shortcut of applying a fast
function to the old constant.  Instead it upgrades the old reflector to a
family oracle, performs autonomous operator reflection, feeds the result back
as the next seed, and repeats the entire strict closure at an input-indexed
number of levels.  The comparison theorem is elementary once the right
invariant---progressivity---is isolated, but the number it certifies is
extraordinarily large.

The most informative exact certificate is not merely $C_{13}<C_{14}$ but
\[
  F_{13}^{44}(42)<F_{14}(42)<C_{14}.
\]
The full construction also dominates all bounded \Fd\ programs using the
packed iterated seed family.  This makes the proposal a substantial
construction-level advance while remaining small enough to audit and integrate.
The remaining mandatory step is a clean Lean~4.34.0 kernel build in the base
repository, using the supplied scripts and audit extensions.

%======================================================================
\appendix
\section{Lean declaration inventory}

\begin{longtable}{>{\ttfamily\raggedright\arraybackslash}p{0.39\textwidth}
                  >{\raggedright\arraybackslash}p{0.53\textwidth}}
\toprule
\normalfont Declaration & Purpose \\
\midrule
\endhead
reflectFdFamily & Family-aware \Fd\ reflector with a finite prefix maximum. \\
fdFamilyTransformer & Base progressive transformer. \\
tree14 & Autonomous schedule stacked above \lean{tree12}. \\
towerStep & Feedback through a full operator-stage closure. \\
operatorTower & Natural-number-indexed strict transformer tower. \\
family13 & Predecessor construction as a function, not a constant. \\
iterateF13 & Right-associated finite iteration of \lean{F13}. \\
seed13 & Family of further \lean{F13} iterations. \\
F14 & Input-indexed tower output. \\
contender\_14 & Final closed natural number. \\
lt\_reflectFdFamily & Generic maximum-to-reflector strict inequality. \\
oracle\_lt\_reflectFdFamily & Visible family coordinate bound. \\
fvalue\_lt\_reflectFdFamily & Bounded \Fd\ program bound. \\
fdFamilyTransformer\_progressive & Base progressivity. \\
transformer\_lt\_towerStep & Strict feedback-step theorem. \\
towerStep\_progressive & Preservation of progressivity. \\
operatorTower\_progressive & Invariant for every tower level. \\
operatorTower\_level\_lt\_succ & Strict pointwise level chain. \\
seed13\_lt\_F14 & Uniform visible-seed lower bound. \\
fortyTwo\_F13\_steps\_lt\_F14 & Exact contest-input margin. \\
contender\_13\_lt\_contender\_14 & Required comparison. \\
fvalue\_lt\_contender\_14 & Outer-budget program domination. \\
\bottomrule
\end{longtable}

\section{Reproduction checklist}

\begin{enumerate}
  \item Check out BusyLean commit
  \lean{1ee50c7d97bbed57a60709f661453c3e213f4db3}.
  \item Confirm \lean{lean-toolchain} contains
  \lean{leanprover/lean4:v4.34.0}.
  \item Run \texttt{scripts/apply\_overlay.sh\ /path/to/BusyLean} from this archive.
  \item Append the supplied audit block to the repository audit.  Its path is
        \path{patches/Audit14.append.lean}.  Alternatively, apply the unified
        patch in \path{patches/integration.patch}.
  \item Run \texttt{python3\ scripts/check.py}.
  \item Run \texttt{lake\ env\ lean\ test/Contender14Audit.lean}.
  \item Save stdout, wall-clock timings, Lean version, operating system, and
        commit hash as the acceptance receipt.
  \item Confirm the dependency-closure assertion excludes all predecessor
        constants and that the printed axiom sets are permitted.
\end{enumerate}

\section{Artifact integrity}
The archive includes \lean{SHA256SUMS.txt}, generated after PDF compilation and
quality inspection.  The file \lean{artifacts/static\_audit.txt} records the
preflight hashes of the two principal Lean modules.  The verification-status
file distinguishes completed checks from checks that still require a pinned
Lean installation.

\begin{thebibliography}{9}
\bibitem{BusyLean}
V.~Reshetnikov,
\emph{BusyLean: name the biggest number, in Lean}, GitHub repository,
2026.  Base commit used here:
\texttt{1ee50c7d97bbed57a60709f661453c3e213f4db3}.

\bibitem{Roux}
C.~Roux,
\emph{name-the-biggest-number}, Coq contest repository.

\bibitem{Aaronson}
S.~Aaronson,
\emph{Who Can Name the Bigger Number?}, expository essay.

\bibitem{Girard}
J.-Y.~Girard,
\emph{Interpr\'etation fonctionnelle et \'elimination des coupures de
l'arithm\'etique d'ordre sup\'erieur}, 1972.

\bibitem{Tait}
W.~W.~Tait,
Intensional interpretations of functionals of finite type I,
\emph{Journal of Symbolic Logic} 32 (1967), 198--212.
\end{thebibliography}

\end{document}
